% ===========================================================================
%  sviluppatore/06-dashboard/07-studio.tex
% ===========================================================================
\section{Chat con l'agente, Genetics e colori della View}\label{sez:studio-pagine}

Tre pagine rifatte con gli stessi criteri: una superficie unica alta quanto
la finestra, classi con un prefisso nuovo, gli id di sempre dove altri script
li cercano, e sul server solo ciò che serviva a mostrare di più (la curva
della fitness, i metadati completi, le preferenze dei colori).

\subsection{Tre regole comuni}

\begin{description}[leftmargin = 1.4em, style = nextline, itemsep = 4pt]
  \item[Prefissi nuovi] \file{ui-system.css} si carica \emph{dopo} il foglio
    della pagina e dà fondo, raggio e bordo con \cod{!important} alle classi
    di prima (\cod{.ag-side}, \cod{.ag-main}, \cod{.ag-chat},
    \cod{.ag-composer}, \cod{.gn-card}, \cod{.gn-form}, \cod{.gn-right},
    \cod{.gn-seg}, \cod{.gn-tbl}\ldots). Le pagine nuove usano
    \cod{ag-layout}/\cod{ag-rail}/\cod{ag-conv} e il prefisso \cod{gx-};
    restano \cod{gn-} solo le classi che scrive \file{genetic\_verdetto.js}.
    I contenitori con pulsanti portano \cod{data-ui-group="none"}: altrimenti
    \file{ui-system.js} li classifica come serie e ne cambia il raggio, o
    marca la scelta di una tavolozza con il fondo pieno.
  \item[Alte quanto la finestra] Lo schema di \file{home.css}: il canvas con
    \cod{data-tab} della pagina prende \cod{calc(100vh - var(--qs-tb-h))}, il
    contenuto è una colonna flessibile con \cod{min-height: 0}, e scorrono le
    aree interne. Genetics lo fa sopra i \num{1100} pixel; sotto, le colonne
    si impilano e scorre la pagina, con il pulsante di avvio fermo in fondo.
  \item[Niente \cod{data-i18n} sui testi nuovi] Il dizionario di
    \file{i18n.js} sostituisce il testo di un elemento con la voce della sua
    chiave: una frase nuova con una chiave vecchia tornerebbe quella di
    prima. I dati dell'utente (nomi di chat e strategie, espressioni) stanno
    in elementi \cod{data-user} o dentro \cod{code}/\cod{pre}, che il
    traduttore salta.
\end{description}

\subsection{La chat con l'agente}

\file{templates/agent.html}, \file{static/js/agent.js},
\file{static/css/agent.css} e \file{static/js/ag\_md.js}, il Markdown delle
risposte. \file{ag\_md.js} prima \emph{sfugge} tutto il testo e poi costruisce
la struttura (titoli, elenchi annidati con le caselle, citazioni, tabelle,
codice); i collegamenti passano solo se sono \cod{http} o \cod{https}. Un
blocco di codice con \cod{file:percorso} mostra il nome del file, uno con la
lingua la evidenzia (Python e JSON, classi \cod{tk-*}). I messaggi
dell'utente usano la forma semplice: righe e codice, nient'altro. Si carica
nel browser (\cod{window.AgMd}) e in \pkg{node}, dove la prova lo esegue; lo
usa anche Genetics per la regola in anteprima.

Le risposte portano in coda due commenti HTML che la pagina legge e il
Markdown non mostra: \cod{<!--meta:\{\ldots\}-->} (durata, file scritti con
righe aggiunte e tolte, operazioni con l'esito, tentativi corretti) diventa
il registro e l'elenco dei file; \cod{<!--interfacce:[\ldots]-->} porta i
riquadri, per esempio il confronto fra stream come grafico e tabella.

\simbolo{agent.\_auto\_title} chiede al modello un titolo breve dopo il
primo messaggio. Un modello piccolo a volte risponde alla domanda invece di
intitolarla: oltre otto parole o sessanta caratteri la risposta non è un
titolo, e si usa l'inizio della domanda, tagliato fra due parole. Toglie
anche \enquote{Title:}, il Markdown e il punto finale. Eliminare una chat
passa per \simbolo{window.showConfirm}: in Electron le finestre di dialogo
native lasciano la pagina senza fuoco.

\subsection{Genetics}

\file{templates/genetic.html}, \file{static/js/genetic.js} (la pagina),
\file{static/js/genetic\_verdetto.js} (il minimo di $n$ e il verdetto del
filtro), \file{static/css/genetic.css}. I due script si parlano con eventi su
\cod{document}, così ognuno resta leggibile da solo:

\begin{center}
\begin{tabular}{@{}L{34mm}L{104mm}@{}}
\toprule
\textbf{Evento} & \textbf{Quando, e chi ascolta} \\
\midrule
\rigaTenue
\cod{qs-gn-config} & La configurazione è arrivata e i campi hanno i valori
  veri (predefiniti o dell'ultimo avvio): il verdetto ricalcola i limiti. \\
\cod{qs-gn-ambito} & Le caselle delle categorie sono state ricostruite. \\
\rigaTenue
\cod{qs-gn-job} & Ogni lettura del lavoro; il dettaglio è il lavoro: il
  verdetto si ridisegna. \\
\cod{qs-gn-installata} & Una strategia respinta è stata installata a mano:
  la pagina aggiorna l'elenco e mostra i collegamenti. \\
\bottomrule
\end{tabular}
\end{center}

\noindent
I controlli segmentati (universo, timeframe) sono radio; il valore vero sta
nei campi nascosti \cod{\#gnScopeMode} e \cod{\#gnTf}, perché il verdetto
legge il primo e ne ascolta il \cod{change}. \simbolo{window.gnQuotaNonValida}
dice alla pagina se il filtro permette di partire: il motivo per cui il
pulsante è spento si scrive sopra il pulsante.

\begin{description}[leftmargin = 1.4em, style = nextline, itemsep = 4pt]
  \item[Il minimo degli asset] Uno per categoria e comunque due
    (\simbolo{minimoDi} in \file{genetic\_verdetto.js}): il server alza a 2
    qualunque valore più basso. Con l'ambito \enquote{un simbolo} il minimo è
    1 da entrambe le parti: prima il server chiedeva comunque due asset
    promossi su uno provato, e quell'ambito non partiva mai.
  \item[Il nome] \cod{\_NAME\_RE} accetta da 3 a 40 caratteri, come dice il
    messaggio; la pagina usa la stessa espressione e segnala anche i nomi già
    presi fra le strategie del sistema.
  \item[Il lavoro] Oltre ai campi di prima, \cod{history}: una voce per
    generazione, \cod{\{g, tr, va\}}, più \cod{ag} e \cod{ap} (varianti usate
    e proposte) dove l'agente è intervenuto --- il motore chiama
    \cod{progress} due volte in quelle generazioni, e la seconda annota la
    voce già scritta. Poi \cod{timeframe}, \cod{pop}, \cod{mutation\_model},
    \cod{installata} dopo un'installazione a mano, e \cod{promotion} anche
    per le promosse.
  \item[I metadati] \file{gp\_meta.json} salva lo storico corto di tutte le
    generazioni (prima: quello lungo del motore, con espressione e regola di
    ciascuna, solo le ultime 60). Un'installazione a mano scrive gli stessi
    campi di una promossa --- metriche, curva, mutazioni, mattoni,
    divisione --- che prima mancavano. \percorsoHTTP{/genetic/api/strategies}
    legge le due forme dello storico e aggiunge \cod{approved} e
    \cod{regola}.
  \item[L'ultimo avvio] \percorsoHTTP{/genetic/api/run} scrive le impostazioni
    in \file{genetic\_last\_run.json} (senza il nome);
    \percorsoHTTP{/genetic/api/config} le restituisce in \cod{last}, e la
    pagina tiene solo ciò che esiste ancora.
  \item[I mattoni] \cod{BUILTIN[k]["family"]} in \file{genetic/primitives.py}
    li raggruppa nella pagina; il motore e l'esportazione leggono solo
    \cod{label} e \cod{lines}.
\end{description}

\noindent
Il grafico della fitness è un SVG scritto a mano e ridisegnato a ogni lettura
e a ogni cambio di larghezza (\cod{ResizeObserver}). $-999$ è la fitness delle
regole che non operano: resta fuori scala, e il mirino la chiama
\enquote{inactive}. La tabella delle strategie ha \cod{table-layout: fixed},
perché il dettaglio aperto sotto una riga non la allarghi, e una
\emph{container query} toglie la data sotto i \num{760} pixel e il drawdown
sotto i \num{620}. \enquote{Run backtest} apre \percorsoHTTP{/run?focus=<nome>}:
\simbolo{evidenziaRichiesta} in \file{runner.js} porta la scheda a vista e la
accende per un momento.

\subsection{I colori delle candele nella View}

\percorsoHTTP{/api/visual/prefs}: \cod{GET} i colori salvati, \cod{POST
\{candles\}} li sostituisce, \cod{null} torna ai predefiniti.
\simbolo{\_colori\_candele} tiene solo \cod{\{c: "\#rrggbb", a: 0..100\}}; un
valore diverso torna al predefinito di quel ruolo invece di arrivare al
grafico. Il file è \file{visual\_view\_prefs.json}, accanto a
\file{projects.json}, e il \cod{POST} riscrive solo la voce \cod{candles}.

\begin{perche}{Sul server e non nel browser}
Il \cod{localStorage} vale per origine, e l'origine cambia con la porta del
server locale (5000, poi 47821--47823): i colori scelti oggi sparirebbero al
prossimo avvio su un'altra porta.
\end{perche}

\noindent
Il modello riceve i colori dalla funzione globale \simbolo{colori\_candele}
(\cod{data-candles} su \cod{\#visual-root}): il grafico nasce già con i
colori giusti. In \file{visual.js} \simbolo{candleColorOpts} li traduce nelle
opzioni di Lightweight Charts --- corpo e bordo delle candele in salita da
\voce{Buy}, stoppini da \voce{Weak buy}, lo stesso in discesa ---,
\simbolo{barColorOpts} dà alle barre OHLC i due colori dei corpi e
\simbolo{volumeColor} gli stessi al volume, al 45\,\%.
\simbolo{ccApplica} ridisegna in un \cod{requestAnimationFrame},
\simbolo{ccSalva} salva \num{350} millisecondi dopo l'ultima modifica.

\subsection{Le prove}

\file{tests/studio\_finto.py} prepara una copia dell'app con cinque ETF
sintetici (duemila barre giornaliere, abbastanza per evoluzioni vere in pochi
secondi), sei chat datate rispetto a oggi e un Ollama finto in un thread: tre
modelli, titoli brevi, varianti valide per le mutazioni guidate, una risposta
fissa per la chat. \file{tests/prova\_studio.py} la usa senza browser:
\file{ag\_md.js} in \pkg{node} (anche l'HTML che non deve passare), il
titolo automatico, l'API dei colori, i controlli di Genetics (nome, minimo,
un simbolo), lavori veri con la curva, l'agente, l'installazione a mano, una
promozione forzata, i metadati vecchi, un lavoro fermato, e gli id che gli
script cercano. \file{tests/e2e\_studio.js} avvia la stessa copia e prova le
tre pagine in Chrome nei due temi.

\begin{codiceInterno}{console}{Le prove delle tre pagine}{}
python tests/prova_studio.py
PLAYWRIGHT_DIR=... CHROME_BIN=/usr/bin/google-chrome node tests/e2e_studio.js
python tests/studio_finto.py --porta 5075 --sessione /tmp/sessione.txt
\end{codiceInterno}

\noindent
L'ultima riga lascia l'app finta accesa per lavorare sull'aspetto delle
pagine; l'indirizzo con il gettone si stampa all'avvio.
